\documentclass[11pt,a4paper]{article}

\usepackage[utf8]{inputenc}
\usepackage[T1]{fontenc}
\usepackage{amsmath}
\usepackage{amssymb}
\usepackage{geometry}
\usepackage{graphicx}
\usepackage{hyperref}
\usepackage{natbib}
\usepackage{booktabs}
\usepackage{multirow}
\usepackage{xcolor}
\usepackage{enumitem}
\usepackage{caption}
\usepackage{algorithm}
\usepackage{algorithmic}

\geometry{margin=1in}

\hypersetup{
    colorlinks=true,
    linkcolor=blue!60!black,
    citecolor=blue!60!black,
    urlcolor=blue!60!black
}

\title{\textbf{Role Adherence: A Metric for Evaluating Role Consistency\\
in Multi-Turn Conversational AI}}

\author{
    Axel Fritz\\
    \texttt{axel.fritz@alquimia.ai}\\
    Alquimia AI
    \and
    Alex Fiorenza\\
    \texttt{alex.fiorenza@alquimia.ai}\\
    Alquimia AI
}

\date{April 2026}

\begin{document}

\maketitle

\begin{abstract}
As conversational AI systems are increasingly deployed in role-constrained
settings --- customer support agents, educational tutors, medical assistants
--- the ability to measure whether a model consistently adheres to its
assigned role becomes a core quality requirement. This paper proposes
\textbf{Role Adherence}, a metric for the Gaussia evaluation framework that
quantifies the degree to which an AI assistant maintains its defined role
throughout a multi-turn conversation. We define adherence constructively,
describe a dual-path scoring strategy conditioned on the availability of
ground-truth data, and discuss design decisions and alternative formulations
considered during development. No empirical validation is presented; this
paper constitutes a design proposal intended to guide the implementation of
the metric within the Gaussia SDK.
\end{abstract}

% ---------------------------------------------------------------------------
\section{Introduction}
% ---------------------------------------------------------------------------

Modern LLM-based assistants are typically deployed with a system prompt that
defines their role: the topics they should cover, the tone they should adopt,
the behaviours they should exhibit, and the boundaries they should respect.
Despite this, models frequently deviate from their assigned roles --- answering
out-of-scope questions, abandoning their defined persona, or failing to exhibit
the behaviour their role requires.

Existing evaluation frameworks either measure role adherence at the level of
individual turns in isolation \citep{zheng2023judging}, or rely entirely on
LLM-based judges, which introduce reproducibility and cost concerns
\citep{deriu2021survey}. This paper proposes a more complete metric that
evaluates adherence across the full conversational context, supports both
structured and unstructured role definitions, and offers a deterministic
scoring path when reference data is available.

% ---------------------------------------------------------------------------
\section{Problem Statement}
% ---------------------------------------------------------------------------

Given a multi-turn conversation between a user and an AI assistant, and a
definition of the role the assistant is expected to play, we want to compute a
scalar score representing how consistently the assistant adhered to that role
throughout the conversation.

Let a conversational session be a sequence of assistant turns
$t_1, t_2, \ldots, t_n$, and let $R$ denote the role definition. The Role
Adherence score is defined as:

\begin{equation}
  \text{RoleAdherence}(R, T) = \frac{1}{n} \sum_{i=1}^{n}
  \text{adhere}(t_i,\; T_{<i},\; R)
  \label{eq:score}
\end{equation}

where $T_{<i}$ is the conversation history up to turn $i$, and
$\text{adhere}(\cdot) \in [0, 1]$ is a per-turn adherence score. The
session score lies in $[0, 1]$, where 1 indicates full adherence across
all turns. The range of $\text{adhere}(\cdot)$ and how it is computed
depend on the scoring path selected (Section~\ref{sec:scoring}).

% ---------------------------------------------------------------------------
\section{Defining Adherence}
\label{sec:adherence}
% ---------------------------------------------------------------------------

\subsection{Constructive vs.\ Restrictive Adherence}

Two competing definitions of adherence are possible.

\textbf{Restrictive adherence} asks: \emph{did the assistant avoid going
outside the scope of its role?} A turn is compliant if it does not address
topics or exhibit behaviours explicitly excluded by the role. This definition
only penalises negative deviations.

\textbf{Constructive adherence} asks: \emph{did the assistant actively behave
as its role requires?} A turn is compliant only if it exhibits the behaviour,
tone, and knowledge a properly role-playing assistant should exhibit. This
definition penalises both negative deviations (scope violations) and positive
failures (failing to act as the role demands).

Consider an assistant assigned the role of a technical support agent. If the
user asks a question within the assistant's domain and the assistant responds
``I don't know'', restrictive adherence would record no violation --- no
out-of-scope content was produced. Constructive adherence would record a
failure, because a role-adherent support agent is expected to answer questions
in its domain.

We adopt the \textbf{constructive definition} as the primary formulation for
this metric, as it provides a more meaningful signal for production quality
assessment. Restrictive adherence remains a valid and simpler formulation for
use cases where the primary concern is containment rather than quality of
execution. Both definitions are implementable within the proposed schema.

\subsection{Conversational Context}

When evaluating whether turn $t_i$ adheres to role $R$, the evaluator receives
the full conversation history $T_{<i} = \{t_1, \ldots, t_{i-1}\}$. A turn
evaluated in isolation may appear off-role but be entirely appropriate given
prior context. Conversely, a turn evaluated in isolation may appear neutral but
represent a pattern of escalating deviation only visible across the
conversational arc.

The alternative --- evaluating each turn independently --- is computationally
cheaper and avoids context-window constraints in very long conversations. We
note this as a valid trade-off for resource-constrained deployments but do not
adopt it as the default behaviour.

% ---------------------------------------------------------------------------
\section{Role Definition Schema}
% ---------------------------------------------------------------------------

\subsection{Placement in the Data Model}

The role definition is a property of the conversational session, not of the
metric. Two sessions with the same conversation content but different assigned
roles represent fundamentally different evaluation scenarios. Accordingly,
\texttt{chatbot\_role} is defined as a field on the \texttt{Dataset} object.

An alternative considered was to pass the role as a parameter at metric
instantiation time, which would allow evaluating the same dataset against
multiple roles. This was discarded on the grounds that a dataset already
carries implicit contextual assumptions: the way a user interacts with a
customer support agent differs from how they interact with a coding assistant,
even if the raw conversation text were identical. Separating the role from the
session would break this coupling and weaken the semantic integrity of the
dataset object.

\subsection{Format: String or Structured}

Role definitions in practice range from a single prose sentence to detailed
specifications with topic lists, tone requirements, and explicit restrictions.
To accommodate both, we propose that \texttt{chatbot\_role} support two
formats:

\textbf{String format} --- free-form natural language description of the role.
Flexible and low-friction for exploratory evaluation.

\textbf{Structured format} --- a \texttt{RoleDefinition} object with explicit
fields: \texttt{description} (required), \texttt{allowed\_topics} (optional
list), \texttt{tone} (optional string), and \texttt{restrictions} (optional
list). This format reduces ambiguity in judge prompts and unlocks
constraint-based deterministic evaluation (Section~\ref{sec:constraints}).

Both formats are valid inputs. The evaluation pipeline adapts accordingly. The
structured format is proposed as the preferred option for production
deployments where reproducibility and auditability are priorities.

% ---------------------------------------------------------------------------
\section{Scoring Methodology}
\label{sec:scoring}
% ---------------------------------------------------------------------------

\subsection{Dual-Path Strategy}

The scoring strategy bifurcates based on the availability of ground-truth
data:

\begin{itemize}
  \item \textbf{Without ground truth}: an LLM-as-a-judge evaluates each
        assistant turn against the role definition and conversation history.
  \item \textbf{With ground truth}: a deterministic metric compares each
        assistant turn against an expected role-adherent reference response.
\end{itemize}

This separation is principled. LLM-based judges are well-suited for
open-ended evaluation where no reference answer exists \citep{zheng2023judging},
but introduce non-determinism, model dependency, and inference cost. When
reference data is available, deterministic metrics offer reproducibility,
auditability, and lower computational cost.

\subsection{LLM-as-a-Judge Path}

When no ground truth is available, each assistant turn $t_i$ is evaluated by
an LLM judge that receives: (1) the role definition $R$; (2) the conversation
history $T_{<i}$; and (3) the turn to evaluate $t_i$. The judge prompt is
constructed to implement the constructive adherence definition: the judge is
asked not only whether the turn avoids scope violations, but whether it
actively reflects the expected behaviour of the defined role.

The judge can be configured to return either a \textbf{continuous score} in
$[0, 1]$ or a \textbf{binary score} in $\{0, 1\}$. Continuous scoring is more
expressive --- a response that mostly adheres but deviates in tone might
receive 0.7 rather than a hard 0 --- but LLMs are generally not well-calibrated
for numeric outputs, meaning the same response may receive different scores
across calls. Binary scoring trades nuance for consistency and reproducibility.
We propose binary scoring as the default for the LLM judge path, with
continuous scoring available as a configurable option for deployments that
prioritise granularity over stability. When \texttt{include\_reason} is
enabled, the judge additionally returns a natural-language justification for
each per-turn score.

\subsection{Deterministic Path}
\label{sec:deterministic}

When a \texttt{ground\_truth\_assistant} response is available for each turn,
the following metrics are proposed as scoring functions. The per-turn score is
continuous in $[0, 1]$ and is averaged across turns to produce the session
score.

\subsubsection{BERTScore}

BERTScore \citep{zhang2020bertscore} computes semantic similarity between two
texts by extracting contextual token embeddings from a pre-trained language
model and computing precision, recall, and F1 over token-level cosine
similarities. Unlike surface-level metrics, it captures paraphrastic
equivalence. We adopt BERTScore F1 as the \textbf{primary deterministic
metric}:

\begin{equation}
  \text{BERTScore}_{F_1} = 2 \cdot \frac{P_{\text{BERT}} \cdot R_{\text{BERT}}}
  {P_{\text{BERT}} + R_{\text{BERT}}}
\end{equation}

Role-adherent responses are not expected to reproduce the exact phrasing of a
reference. BERTScore tolerates natural linguistic variation while remaining
independent of a secondary judge model.

\subsubsection{Semantic Cosine Similarity}

Cosine similarity between sentence-level embeddings \citep{reimers2019sentencebert}
is a lighter-weight alternative. A sentence encoder maps each response to a
dense vector; adherence is scored as the cosine of the angle between the
candidate and reference vectors:

\begin{equation}
  \text{sim}(\mathbf{a}, \mathbf{b}) =
  \frac{\mathbf{a} \cdot \mathbf{b}}{\|\mathbf{a}\|\,\|\mathbf{b}\|}
\end{equation}

Cosine similarity operates at sentence granularity rather than token
granularity, making it faster but less sensitive to fine-grained semantic
differences. BERTScore and cosine similarity are compatible and non-exclusive
--- they measure related but distinct levels of semantic granularity. Where
computational resources permit, BERTScore is preferred.

\subsubsection{Natural Language Inference}

NLI models \citep{bowman2015snli} classify the relationship between two texts
as \emph{entailment}, \emph{neutral}, or \emph{contradiction}. Applied to role
adherence, a candidate response that contradicts what a ground-truth
role-adherent response asserts is a strong signal of non-adherence. NLI is
proposed as a complementary signal for detecting explicit violations rather
than measuring overall similarity. A \emph{contradiction} label serves as a
hard penalty regardless of similarity scores.

\subsubsection{Constraint-Based Evaluation}
\label{sec:constraints}

When the role is defined in structured format, adherence can be partially
evaluated through explicit constraint verification without any model inference:

\begin{itemize}
  \item \textbf{Topic scope}: whether the response discusses topics outside
        \texttt{allowed\_topics}.
  \item \textbf{Restriction compliance}: whether the response violates any
        item in \texttt{restrictions}.
  \item \textbf{Tone consistency}: whether the detected tone matches the
        expected \texttt{tone} field.
\end{itemize}

These checks are deterministic, fully reproducible, and computationally
trivial. They are proposed as a first-pass filter that complements
BERTScore or cosine similarity rather than replacing them.

\subsubsection{Lexical Metrics: Considered and Discarded}

ROUGE \citep{lin2004rouge} and BLEU \citep{papineni2002bleu} measure n-gram
overlap between candidate and reference texts. Both were considered and
discarded as primary metrics for this use case.

Role-adherent responses are not expected to reproduce the exact phrasing of a
reference --- they are expected to exhibit the same role-consistent behaviour
with natural linguistic variation. N-gram metrics penalise paraphrases and
reward surface-level copying, neither of which is desirable here. ROUGE and
BLEU are noted as valid lightweight baselines for benchmarking purposes only.

% ---------------------------------------------------------------------------
\section{Configuration Parameters}
% ---------------------------------------------------------------------------

The metric exposes four optional parameters that control evaluation behaviour
without altering the underlying methodology.

\textbf{\texttt{threshold}} (default: 0.5) --- the minimum session score
required to consider the assistant as role-adherent. Scores at or above this
value are marked as passing. The threshold applies to the final session scalar
produced by Equation~\ref{eq:score}.

\textbf{\texttt{strict\_mode}} (default: \texttt{False}) --- when set to
\texttt{True}, the session score is converted to a binary output: 1 if the
score equals 1.0 (every turn adheres), 0 otherwise. This also overrides
\texttt{threshold} to 1.0. Strict mode is intended for high-criticality
deployments --- such as medical or legal assistants --- where partial adherence
is not acceptable.

\textbf{\texttt{include\_reason}} (default: \texttt{True}) --- when enabled
and the LLM-judge path is active, the judge returns a natural-language
justification alongside each per-turn score. Justifications aid interpretability
and are included in the metric output for inspection.

\textbf{\texttt{model}} --- the LLM used as the judge in the no-ground-truth
path. Follows the \texttt{Guardian} abstraction already present in the Gaussia
framework, allowing any supported provider to be substituted without changes
to the metric logic.

% ---------------------------------------------------------------------------
\section{Implementation Considerations}
% ---------------------------------------------------------------------------

The metric extends the \texttt{Gaussia} base class. The \texttt{batch()}
method iterates over each assistant turn in a session, selects the
LLM-judge or deterministic path based on the presence of
\texttt{ground\_truth\_assistant}, and appends a per-turn score. The session
score is the mean of per-turn scores (Equation~\ref{eq:score}).

The \texttt{chatbot\_role} field is read from the \texttt{Dataset} object.
When \texttt{chatbot\_role} is a \texttt{RoleDefinition}, the structured
fields are serialised into the judge prompt and, where applicable, used for
constraint-based evaluation. When it is a plain string, it is passed directly.

Statistical mode compatibility (Frequentist and Bayesian) follows the existing
Gaussia pattern. No changes to the statistical layer are required.

\newpage

% ---------------------------------------------------------------------------
\section{Design Decisions Summary}
% ---------------------------------------------------------------------------

\begin{table}[h]
  \centering
  \caption{Summary of design decisions and alternatives considered.}
  \label{tab:decisions}
  \begin{tabular}{p{4.2cm} p{5.2cm} p{4.5cm}}
    \toprule
    \textbf{Dimension} & \textbf{Decision} & \textbf{Alternative Considered} \\
    \midrule
    \texttt{chatbot\_role} location
      & Field on \texttt{Dataset}
      & Parameter of the metric \\
    \addlinespace
    Adherence definition
      & Constructive
      & Restrictive \\
    \addlinespace
    Role format
      & \texttt{str | RoleDefinition}
      & String-only \\
    \addlinespace
    Judge context
      & Full conversation history up to turn
      & Per-turn isolation \\
    \addlinespace
    Scoring without GT
      & LLM-as-a-judge
      & --- \\
    \addlinespace
    LLM judge output mode
      & Binary (default), continuous (configurable)
      & Continuous-only \\
    \addlinespace
    Scoring with GT
      & BERTScore (primary), cosine similarity (lightweight), NLI (violation detection)
      & ROUGE, BLEU (discarded) \\
    \addlinespace
    Score output
      & Session-level scalar in $[0, 1]$
      & --- \\
    \bottomrule
  \end{tabular}
\end{table}


% ---------------------------------------------------------------------------
\section{Future Work}
% ---------------------------------------------------------------------------

The current proposal scopes Role Adherence to text-only conversations.
Extending the metric to \textbf{multimodal settings} is identified as the
primary direction for future development.

\subsection{Multimodal Role Adherence}

Conversational AI systems are increasingly capable of processing and generating
content across multiple modalities --- images, audio, and video alongside text.
A role-adherent response in a multimodal context must satisfy the same
constructive definition proposed in Section~\ref{sec:adherence}: the assistant
must actively behave as the role requires, not merely avoid out-of-scope
output. This requirement extends naturally to non-text modalities.

\subsubsection{Role Definition Extensions}

The \texttt{RoleDefinition} schema would require a \texttt{modalities} field
specifying which content types the role operates over (e.g.,
\texttt{["text", "image"]}, \texttt{["text", "audio"]}). The presence of this
field makes modality compliance an explicit dimension of adherence: an
assistant whose role requires generating images but responds with text only
fails constructive adherence regardless of the textual content quality.

\subsubsection{Scoring in Multimodal Settings}

The dual-path strategy generalises to multimodal content, but the concrete
scoring functions differ by modality:

\textbf{LLM-as-a-judge path} --- when the conversation contains non-text
content, the judge must be a vision-capable or otherwise multimodal model.
The judge logic and prompt structure remain unchanged; only the model
requirement is elevated.

\textbf{Deterministic path} --- BERTScore and cosine similarity over text
embeddings do not apply to image or audio content. For image modalities, we
propose exploring CLIP-based similarity \citep{radford2021clip}, which embeds
images and text into a shared semantic space and allows cross-modal comparison.
For mixed responses combining text and image, a weighted combination of
text-based and image-based similarity scores would be required, with the
weighting strategy to be determined empirically.

\subsubsection{Data Schema}

The \texttt{Batch} schema currently models \texttt{query} and
\texttt{assistant} as plain text fields. Supporting multimodal content would
require these fields to accept structured payloads that can carry text, image
references, or audio alongside appropriate metadata. This is a schema-level
change with no impact on the metric formula or scoring logic.

\subsection{Additional Directions}

\begin{itemize}
  \item \textbf{Empirical validation}: comparison of scoring strategies
        (BERTScore, cosine similarity, NLI, LLM judge) against human
        annotations of role adherence across diverse role types and domains.

  \item \textbf{Structured role extraction}: automatic construction of a
        \texttt{RoleDefinition} from a free-form system prompt using an LLM,
        enabling constraint-based evaluation without manual annotation.

  \item \textbf{Turn-level reporting}: exposing per-turn adherence scores
        alongside the session aggregate to allow fine-grained diagnosis of
        where role deviation occurs.

  \item \textbf{Cross-lingual evaluation}: role adherence in multilingual
        assistants, where role definitions and responses may be in different
        languages.
\end{itemize}

% ---------------------------------------------------------------------------
\section{Conclusion}
% ---------------------------------------------------------------------------

We have proposed Role Adherence, a metric for the Gaussia evaluation framework
that measures how consistently an AI assistant maintains its assigned role
across a multi-turn conversation. The metric adopts a constructive definition
of adherence, evaluates each turn in full conversational context, and offers a
dual scoring path --- LLM-as-a-judge when no reference data is available, and
deterministic semantic metrics (BERTScore, cosine similarity, NLI) when ground
truth is present. The role definition is modelled as a first-class field on the
session dataset, supporting both free-form string and structured formats. This
paper constitutes a design proposal scoped to text-only conversations; extension
to multimodal settings is identified as the primary direction for future work.

\bibliographystyle{plainnat}
\bibliography{refs}

\end{document}
